%% =====================================================================
%%  CAPITULO 7 - IMPLEMENTACAO
%% =====================================================================

\chapter{Implementação}
\label{cap:implementacao}

Este capítulo descreve a implementação do sistema-alvo: a estrutura do artefato,
a máquina de estados da partida, o laço de animação, o algoritmo de um passo de
simulação, a fila de entrada, o sorteio do alvo, a persistência, a renderização e
as decisões de acessibilidade. Cada seção aponta para a decisão arquitetural e o
requisito correspondentes.

\section{Estrutura do artefato}
\label{sec:impl-estrutura}

O produto é um único arquivo de 1\,063 linhas e 32\,305 bytes, organizado em
treze seções comentadas. A Figura~\ref{fig:artefato} representa sua estrutura.

%% ---------------------------------------------------------------------
%%  FIGURA - Estrutura do artefato unico (DG-07)
%% ---------------------------------------------------------------------
\begin{figure}[htb]
	\caption{Estrutura do artefato único e suas treze seções de código}
	\label{fig:artefato}
	\centering
	\begin{tikzpicture}[
			font=\footnotesize,
			bloco/.style={draw, rounded corners=2pt, align=center, fill=black!4, inner sep=4pt},
			interno/.style={draw, rounded corners=1pt, align=left, fill=black!8, font=\scriptsize, inner sep=3pt},
			>=Stealth,
			fl/.style={->, thick}
		]

		\node[bloco, minimum width=3.1cm, minimum height=0.8cm] (head) at (0,0) {\texttt{head}\\ metadados e título};
		\node[bloco, minimum width=3.1cm, minimum height=0.8cm, below=0.3cm of head] (style) {\texttt{style}\\ tema e layout};
		\node[bloco, minimum width=3.1cm, minimum height=0.8cm, below=0.3cm of style] (body) {\texttt{body}\\ HUD, tela, D-Pad};
		\node[bloco, minimum width=3.1cm, minimum height=0.8cm, below=0.3cm of body] (script) {\texttt{script}\\ lógica do jogo};

		\node[interno, right=0.9cm of script, anchor=west] (sec) {%
			1 configuração \hfill 8 passo de simulação\\
			2 referências de DOM \hfill 9 renderização\\
			3 estado \hfill 10 laço de animação\\
			4 utilitários \hfill 11 interface\\
			5 persistência \hfill 12 entrada\\
			6 regras puras \hfill 13 inicialização\\
			7 fluxo de partida \hfill \\
		};

		\draw[fl] (script) -- (sec);

		\node[font=\scriptsize, above=0.25cm of head] (tag) {Arquivo único: 1\,063 linhas, 32\,305 bytes};
	\end{tikzpicture}
	\legend{Fonte: elaborado pelo autor (2026) a partir de \texttt{index.html}. A restrição de artefato
	único (\texttt{CON-001}) impede a separação física entre módulos, de modo que o isolamento é obtido
	por funções de responsabilidade única dentro do mesmo arquivo.}
\end{figure}


A organização em seções numeradas é a resposta a \texttt{NFR-005} (código
comentado e organizado) e cumpre função operacional, não apenas estética: como o
processo de verificação precisa localizar trechos específicos por inspeção
textual, a previsibilidade da estrutura reduz o esforço de verificação.

A restrição \texttt{CON-001} (artefato único) tem uma consequência de projeto que
vale explicitar: sem módulos separados, não há como isolar a lógica de simulação
da lógica de apresentação por meio de fronteiras de arquivo. O isolamento foi
obtido por meio de funções com responsabilidade única dentro do mesmo arquivo, e
a separação é garantida por convenção, não por mecanismo. É um enfraquecimento
real da arquitetura em relação a uma solução modular --- aceito porque exigido
pela especificação.

\section{Máquina de estados da partida}
\label{sec:impl-fsm}

A partida é modelada como máquina de estados finitos com cinco estados:
\texttt{ready}, \texttt{running}, \texttt{paused}, \texttt{over} e \texttt{won}.
A Figura~\ref{fig:fsm} apresenta as transições.

%% ---------------------------------------------------------------------
%%  FIGURA - Maquina de estados da partida (DG-04)
%% ---------------------------------------------------------------------
\begin{figure}[htb]
	\caption{Máquina de estados da partida e transições}
	\label{fig:fsm}
	\centering
	\begin{tikzpicture}[
			font=\footnotesize,
			estado/.style={draw, rounded corners=3pt, minimum width=2.4cm, minimum height=0.85cm, align=center},
			inicial/.style={estado, very thick, fill=black!8},
			>=Stealth,
			tr/.style={->, thick},
			rot/.style={font=\scriptsize, inner sep=2pt, fill=white}
		]

		\node[inicial] (ready) at (0,0) {\texttt{ready}};
		\node[estado] (run) at (4.6,0) {\texttt{running}};
		\node[estado] (pause) at (4.6,-2.3) {\texttt{paused}};
		\node[estado] (over) at (9.2,1.1) {\texttt{over}};
		\node[estado] (won) at (9.2,-1.1) {\texttt{won}};

		\draw[tr] (ready) -- node[rot, above] {direção válida} (run);
		\draw[tr] (run) -- node[rot, right] {pausa ou perda de foco} (pause);
		\draw[tr] (pause) -- node[rot, left] {retomada} (run);
		\draw[tr] (run) -- node[rot, pos=0.62, above] {colisão} (over);
		\draw[tr] (run) -- node[rot, pos=0.62, below] {grade completa} (won);

		\draw[tr] (run) to[out=110, in=70, looseness=7] node[rot, above] {passo de simulação} (run);

		\draw[tr] (over) to[out=200, in=30, looseness=1.3] node[rot, pos=0.42, above] {reinício} (ready);
		\draw[tr] (won) to[out=200, in=-10, looseness=1.15] node[rot, pos=0.4, below] {reinício} (ready);
	\end{tikzpicture}
	\legend{Fonte: elaborado pelo autor (2026) a partir de \texttt{index.html}, seção~7.
	O estado \texttt{won} não constava da solicitação original: ele foi criado para tratar o caso de
	borda ``grade completamente ocupada''.}
\end{figure}


Duas decisões de modelagem merecem registro. A primeira é a existência do estado
\texttt{won}, que não constava da solicitação original. Ele é necessário porque a
especificação declarou o caso de borda ``grade completamente ocupada''
(Seção~\ref{sec:esp-bordas}): sem esse estado, o sorteio do alvo retornaria um
valor inválido e o jogo entraria em estado inconsistente. A modelagem do caso de
borda forçou a criação de um estado. Esse é um exemplo concreto de como declarar
casos de borda altera o projeto, e não apenas os testes.

A segunda é que o estado \texttt{ready} não é apenas decorativo. Ele existe
porque a especificação exige que a movimentação comece somente após a primeira
entrada direcional válida (\texttt{FR-012}). A consequência é que a primeira
entrada precisa ser tratada de forma diferente das demais: como o jogador ainda
não se moveu, qualquer direção cardinal é aceita, inclusive a oposta à direção
inicial, e o corpo é reposicionado atrás da nova cabeça para evitar autocolisão
imediata. Sem esse tratamento, pressionar ``esquerda'' na tela inicial mataria o
jogador no primeiro passo.

\section{Laço de animação}
\label{sec:impl-loop}

O laço usa o mecanismo de requisição de quadro de animação
(\citeonline{mdn2024raf}), com acumulador de tempo para separar a simulação da
apresentação, conforme \texttt{FR-014}, \texttt{CON-004} e ADR-002. A
Figura~\ref{fig:game-loop} representa o fluxo.

%% ---------------------------------------------------------------------
%%  FIGURA - Laco de animacao com acumulador (DG-05)
%% ---------------------------------------------------------------------
\begin{figure}[htb]
	\caption{Laço de animação com acumulador de tempo e execução de passos}
	\label{fig:game-loop}
	\centering
	\begin{tikzpicture}[
			font=\footnotesize,
			passo/.style={draw, rounded corners=2pt, minimum width=3.5cm, minimum height=0.8cm, align=center, fill=black!4},
			dec/.style={draw, diamond, aspect=2.4, inner sep=1pt, align=center, font=\scriptsize},
			>=Stealth,
			fl/.style={->, thick},
			rot/.style={font=\scriptsize, inner sep=2pt}
		]

		\node[passo] (raf) at (0,0) {Requisição de quadro};
		\node[passo, below=0.55cm of raf] (delta) {Calcular intervalo $\Delta$ e limitar a 250\,ms};
		\node[dec, below=0.75cm of delta] (rodando) {estado =\\ \texttt{running}?};
		\node[passo, right=2.5cm of rodando] (acum) {Acumular: $a \mathrel{+}= \Delta$};
		\node[dec, below=0.85cm of acum] (suf) {$a \geq t$ e\\ passos $< 5$?};
		\node[passo, right=2.3cm of suf] (tick) {Executar passo\\ de simulação};
		\node[passo, below=0.7cm of tick] (dec2) {Subtrair $t$ de $a$\\ e incrementar contador};
		\node[passo, below=2.2cm of rodando] (desc) {Descartar acúmulo residual};
		\node[passo, below=0.7cm of desc] (render) {Renderizar};
		\node[passo, below=0.7cm of render] (loop) {Agendar próximo quadro};

		\draw[fl] (raf) -- (delta);
		\draw[fl] (delta) -- (rodando);
		\draw[fl] (rodando) -- node[rot, above] {sim} (acum);
		\draw[fl] (acum) -- (suf);
		\draw[fl] (suf) -- node[rot, above] {sim} (tick);
		\draw[fl] (tick) -- (dec2);
		\draw[fl] (dec2) |- (acum);
		\draw[fl] (suf) |- node[rot, pos=0.75, left] {não} (desc);
		\draw[fl] (rodando) -- node[rot, left] {não} (desc);
		\draw[fl] (desc) -- (render);
		\draw[fl] (render) -- (loop);
		\draw[fl] (loop) -- ++(2.2,0) |- (raf);
	\end{tikzpicture}
	\legend{Fonte: elaborado pelo autor (2026) a partir de \texttt{index.html}, seção~10.
	O teto de cinco passos por quadro e o descarte do acúmulo residual são as duas defesas contra a
	espiral da morte.}
\end{figure}


O trecho central está no Quadro~\ref{qua:codigo-loop}.

\begin{quadro}[htb]
	\caption{Trecho do laço de animação: acúmulo de tempo e execução de passos (excertado)}
	\label{qua:codigo-loop}
	\begin{lstlisting}[style=codigo]
const delta = Math.min(timestamp - state.lastTs, 250);
state.lastTs = timestamp;
state.elapsed += delta;

if (state.status === 'running') {
  state.acc += delta;
  let steps = 0;
  while (state.acc >= state.stepMs && steps < CONFIG.maxTicksPerFrame) {
    state.acc -= state.stepMs;
    steps += 1;
    tick();
    if (state.status !== 'running') { state.acc = 0; break; }
  }
  if (state.acc > state.stepMs * CONFIG.maxTicksPerFrame) state.acc = 0;
} else {
  state.acc = 0;
}
	\end{lstlisting}
	\legend{Fonte: \texttt{index.html}, seção~10.}
\end{quadro}

Três detalhes do trecho são decisões de engenharia, não acidentes de codificação.

\textbf{Limitação do passo de tempo.} A primeira linha limita o intervalo a
250\,ms. Sem esse limite, um retorno de segundo plano produziria um intervalo de
vários segundos e o jogador avançaria dezenas de células em um único quadro, sem
possibilidade de reação. O limite converte uma anomalia de temporização em uma
perda controlada de fluidez.

\textbf{Teto de passos por quadro.} A condição \texttt{steps < maxTicksPerFrame}
limita a quantidade de simulações executadas em um mesmo quadro. Sem esse teto, a
recuperação do atraso acumulado consumiria mais tempo de processamento do que o
atraso original, ampliando a defasagem --- fenômeno conhecido como espiral da
morte. O teto rompe o ciclo ao custo de a simulação ficar momentaneamente atrás
do tempo real.

\textbf{Descarte do acúmulo residual.} A linha final descarta o acúmulo quando
ele excede o teto multiplicado pelo intervalo. É a segunda linha de defesa contra
a mesma anomalia.

A Tabela~\ref{tab:tempos} registra as grandezas temporais do laço.

\begin{table}[htb]
	\caption{Grandezas temporais do laço de animação}
	\label{tab:tempos}
	\centering
	\footnotesize
	\begin{tabularx}{\textwidth}{@{}p{4.4cm} c X@{}}
		\toprule
		\textbf{Grandeza} & \textbf{Valor} & \textbf{Observação} \\
		\midrule
		Intervalo inicial entre passos & 140\,ms & Aproximadamente sete passos por segundo \\
		\addlinespace
		Intervalo mínimo & 68\,ms & Aproximadamente quinze passos por segundo \\
		\addlinespace
		Teto de passo de tempo & 250\,ms & Limite do intervalo aceito em um quadro \\
		\addlinespace
		Teto de passos por quadro & 5 & Proteção contra espiral da morte \\
		\addlinespace
		Taxa de quadro alvo & --- & Delegada ao navegador pelo mecanismo de requisição de quadro \\
		\bottomrule
	\end{tabularx}
	\legend{Fonte: elaborado pelo autor (2026) a partir de \texttt{index.html}.}
\end{table}

\section{Algoritmo de um passo de simulação}
\label{sec:impl-tick}

O passo de simulação é a unidade atômica do jogo. Sua ordem de operações é
determinada por ADR-004 e está representada na Figura~\ref{fig:tick}.

%% ---------------------------------------------------------------------
%%  FIGURA - Algoritmo de um passo de simulacao (DG-06)
%% ---------------------------------------------------------------------
\begin{figure}[htb]
	\caption{Algoritmo de um passo de simulação, com destaque para a ordem de ADR-004}
	\label{fig:tick}
	\centering
	\begin{tikzpicture}[
			font=\footnotesize,
			acao/.style={draw, rounded corners=2pt, minimum width=6.4cm, minimum height=0.78cm, align=center, fill=black!4},
			critico/.style={acao, fill=black!14, very thick},
			saida/.style={draw, rounded corners=2pt, minimum width=3.2cm, minimum height=0.78cm, align=center, fill=black!8},
			dec/.style={draw, diamond, aspect=3.0, inner sep=1pt, align=center, font=\scriptsize},
			>=Stealth,
			fl/.style={->, thick},
			rot/.style={font=\scriptsize, inner sep=2pt, fill=white}
		]

		\node[acao] (a1) at (0,0) {1. Aplicar a próxima direção válida da fila};
		\node[acao, below=0.42cm of a1] (a2) {2. Calcular a nova cabeça: cabeça $+$ direção};
		\node[dec, below=0.6cm of a2] (d1) {cabeça fora\\ da grade?};
		\node[acao, below=0.6cm of d1] (a3) {3. Verificar se a cabeça alcançou o alvo};
		\node[critico, below=0.42cm of a3] (a4) {4. Se não cresceu: remover a cauda};
		\node[dec, below=0.6cm of a4] (d2) {cabeça em célula\\ do corpo?};
		\node[acao, below=0.6cm of d2] (a5) {5. Inserir a nova cabeça no início do corpo};
		\node[acao, below=0.42cm of a5] (a6) {6. Se cresceu: somar pontos, reduzir o intervalo, sortear novo alvo};

		\node[saida, right=1.5cm of d1] (o1) {Encerrar partida};
		\node[saida, right=1.5cm of d2] (o2) {Encerrar partida};
		\node[saida, right=1.5cm of a6, align=center] (o3) {Ganhar\\ (sem célula livre)};

		\draw[fl] (a1) -- (a2);
		\draw[fl] (a2) -- (d1);
		\draw[fl] (d1) -- node[rot, left] {não} (a3);
		\draw[fl] (d1) -- node[rot, above] {sim} (o1);
		\draw[fl] (a3) -- (a4);
		\draw[fl] (a4) -- (d2);
		\draw[fl] (d2) -- node[rot, left] {não} (a5);
		\draw[fl] (d2) -- node[rot, above] {sim} (o2);
		\draw[fl] (a5) -- (a6);
		\draw[fl] (a6) -- (o3);
	\end{tikzpicture}
	\legend{Fonte: elaborado pelo autor (2026) a partir de \texttt{index.html}, seção~8.
	O passo 4 é destacado porque sua posição \emph{antes} do teste de colisão (ADR-004) determina se
	ocupar a célula recém-liberada pela cauda é sobrevivência ou morte.}
\end{figure}


O Quadro~\ref{qua:codigo-tick} apresenta o trecho correspondente.

\begin{quadro}[htb]
	\caption{Trecho do passo de simulação: cálculo da cabeça, colisão e remoção da cauda (excertado)}
	\label{qua:codigo-tick}
	\begin{lstlisting}[style=codigo]
if (state.queue.length > 0) {
  state.dir = state.queue.shift();
}
const head = {
  x: state.snake[0].x + state.dir.x,
  y: state.snake[0].y + state.dir.y
};
if (head.x < 0 || head.y < 0 ||
    head.x >= CONFIG.cols || head.y >= CONFIG.rows) {
  endGame('collision');
  return;
}
const growing = state.food !== null && sameCell(head, state.food);
if (!growing) {
  state.snake.pop();
}
	\end{lstlisting}
	\legend{Fonte: \texttt{index.html}, seção~8.}
\end{quadro}

A ordem tem três propriedades que importam.

\textbf{A direção é aplicada antes do movimento.} A fila é consumida no início do
passo, de modo que a direção usada no cálculo corresponde exatamente à intenção
vigente naquele passo.

\textbf{A colisão de borda é testada antes de qualquer mutação do corpo.} Se a
cabeça saiu da grade, o estado não foi alterado. Isso mantém a consistência do
estado observável em caso de encerramento.

\textbf{A cauda é removida antes do teste de colisão com o corpo.} Essa é a
decisão ADR-004 e produz o comportamento declarado no caso de borda~5: mover para
a célula que a cauda acabou de liberar é sobrevivência, não morte. A ordem inversa
produziria a morte.

\section{Fila de entrada e bloqueio de inversão}
\label{sec:impl-entrada}

O bloqueio de inversão instantânea (\texttt{FR-005}) é implementado por fila de
capacidade limitada com validação contra a última direção de referência. A
Figura~\ref{fig:sequencia} apresenta a sequência de tratamento de uma entrada.

%% ---------------------------------------------------------------------
%%  FIGURA - Sequencia de tratamento de entrada (DG-08)
%% ---------------------------------------------------------------------
\begin{figure}[htb]
	\caption{Sequência de tratamento de uma entrada direcional}
	\label{fig:sequencia}
	\centering
	\begin{tikzpicture}[
			font=\scriptsize,
			ator/.style={draw, rounded corners=2pt, minimum width=2.3cm, minimum height=0.7cm, align=center, fill=black!8},
			>=Stealth,
			msg/.style={->, thick},
			ret/.style={->, thick, dashed}
		]

		% atores e linhas de vida
		\node[ator] (jog) at (0,0) {Jogador};
		\node[ator] (ui) at (3.4,0) {Interface\\ (teclado / D-Pad)};
		\node[ator] (fila) at (7.0,0) {Fila de\\ entrada};
		\node[ator] (val) at (10.4,0) {Validação de\\ direção};
		\node[ator] (pas) at (13.9,0) {Passo de\\ simulação};

		\foreach \x in {0,3.4,7.0,10.4,13.9} {
				\draw[dotted, thick] (\x,-0.45) -- (\x,-7.2);
			}

		\draw[msg] (0,-1.0) -- node[above, font=\scriptsize] {tecla ou toque} (3.4,-1.0);
		\draw[msg] (3.4,-1.8) -- node[above, font=\scriptsize] {nome da direção} (7.0,-1.8);
		\draw[msg] (7.0,-2.6) -- node[above, font=\scriptsize] {direção solicitada} (10.4,-2.6);

		\draw[msg] (10.4,-3.4) -- ++(1.0,0) -- ++(0,-0.7) -- node[right, font=\scriptsize, pos=0.55] {descartar se igual} (11.4,-4.1);
		\draw[msg] (10.4,-4.6) -- ++(1.0,0) -- ++(0,-0.7) -- node[right, font=\scriptsize, pos=0.55] {descartar se oposta} (11.4,-5.3);

		\draw[msg] (10.4,-6.0) -- node[above, font=\scriptsize] {enfileirar} (13.9,-6.0);
		\draw[ret] (13.9,-6.7) -- node[above, font=\scriptsize] {aplicar no próximo passo} (7.0,-6.7);

		\node[font=\scriptsize, align=left, anchor=west] at (0,-7.8)
		{Referência de comparação: a última direção enfileirada, ou a direção corrente quando a fila está vazia.};
	\end{tikzpicture}
	\legend{Fonte: elaborado pelo autor (2026) a partir de \texttt{index.html}, seção~6 (ADR-003).
	Comparar com a última direção \emph{enfileirada} permite duas mudanças de noventa graus no mesmo passo
	sem abrir a possibilidade de inversão de cento e oitenta graus.}
\end{figure}


O Quadro~\ref{qua:codigo-fila} apresenta o trecho.

\begin{quadro}[htb]
	\caption{Trecho do enfileiramento de direção com bloqueio de inversão (excertado)}
	\label{qua:codigo-fila}
	\begin{lstlisting}[style=codigo]
const reference = state.queue.length > 0
  ? state.queue[state.queue.length - 1]
  : state.dir;

if (sameCell(next, reference) || isOpposite(next, reference)) {
  return;
}
if (state.queue.length >= CONFIG.maxQueue) {
  state.queue.shift();
}
state.queue.push({ ...next });
	\end{lstlisting}
	\legend{Fonte: \texttt{index.html}, seção~6.}
\end{quadro}

A validação usa como referência a \emph{última direção enfileirada}, e não a
direção corrente, quando a fila não está vazia. A diferença é o que permite duas
mudanças de noventa graus em um mesmo passo curto --- por exemplo, direita seguida
de baixo --- sem permitir a inversão de cento e oitenta graus. Comparar apenas
com a direção corrente rejeitaria a segunda entrada, e o jogador perderia o
comando; comparar com a fila preserva a intenção sem abrir a possibilidade de
inversão.

\section{Sorteio do alvo}
\label{sec:impl-alvo}

O sorteio do alvo (\texttt{FR-006}) impede que o alvo apareça sobre o corpo. A
estratégia tem duas fases. Na primeira, sorteia-se uma célula aleatória e verifica-se
sua presença em um conjunto de ocupação, por até 500 tentativas. Se todas
falharem, executa-se uma varredura determinística em ordem de linha e coluna,
retornando a primeira célula livre.

A escolha de duas fases é uma decisão de robustez. A amostragem aleatória é
adequada enquanto a ocupação é baixa; quando o tabuleiro fica quase cheio, a
probabilidade de acerto cai e o custo esperado de tentativas cresce. A varredura
determinística garante término e, se nenhuma célula estiver livre, retorna valor
nulo, que o passo de simulação interpreta como vitória --- o estado
\texttt{won}.

O Quadro~\ref{qua:codigo-alvo} apresenta o núcleo do sorteio.

\begin{quadro}[htb]
	\caption{Trecho do sorteio de célula livre para o alvo (excertado)}
	\label{qua:codigo-alvo}
	\begin{lstlisting}[style=codigo]
const occupied = new Set(state.snake.map(function (cell) {
  return cell.x + ',' + cell.y;
}));

for (let attempt = 0; attempt < CONFIG.maxFoodAttempts; attempt += 1) {
  const cell = { x: randomInt(CONFIG.cols), y: randomInt(CONFIG.rows) };
  if (!occupied.has(cell.x + ',' + cell.y)) {
    return cell;
  }
}
	\end{lstlisting}
	\legend{Fonte: \texttt{index.html}, seção~6.}
\end{quadro}

Uma nota de eficiência: a codificação da célula como texto para uso em conjunto
tem custo de alocação. A alternativa sem alocação seria uma matriz de ocupação
reutilizada entre passos, o que exigiria limpeza seletiva a cada passo e
introduziria risco de estado residual. A escolha pela legibilidade e
previsibilidade, com o custo de alocação assumido, é coerente com
\texttt{NFR-005} e com o tamanho da grade (24 por 24 células).

\section{Persistência do recorde}
\label{sec:impl-persistencia}

A persistência (\texttt{FR-010}) está isolada em duas funções envoltas em
tratamento de exceção, conforme ADR-006. A gravação ocorre no momento em que a
pontuação supera o recorde, e não apenas no encerramento da partida.

A escolha de gravar durante a partida tem uma consequência favorável: se a página
for fechada abruptamente, o recorde conquistado não se perde. O custo é uma
escrita no armazenamento por item consumido, custo irrelevante na escala do jogo.
A leitura ocorre uma única vez, na inicialização.

O tratamento de exceção cobre dois cenários reais: navegação privada com
armazenamento restrito e política de armazenamento desabilitada pelo usuário. Em
ambos, a função de leitura retorna zero e o jogo permanece jogável com o recorde
válido apenas na sessão corrente. A falha é degradada, não propagada.

\section{Renderização}
\label{sec:impl-render}

A renderização ocorre em superfície vetorial bidimensional
(\citeonline{mdn2024canvas}) e segue ADR-001. Três decisões a caracterizam.

\textbf{Ajuste à densidade da tela.} O Quadro~\ref{qua:codigo-resize} mostra o
ajuste, que atende \texttt{NFR-006}. A razão de pixels é limitada a~2 para evitar
custo de desenho desnecessário em telas de densidade muito alta, nas quais o
ganho visual seria imperceptível e o custo de memória, não.

\begin{quadro}[htb]
	\caption{Trecho do ajuste da superfície de desenho à densidade da tela (excertado)}
	\label{qua:codigo-resize}
	\begin{lstlisting}[style=codigo]
const rect = dom.canvas.getBoundingClientRect();
const cssSize = Math.max(1, Math.round(rect.width));
const dpr = Math.min(window.devicePixelRatio || 1, 2);

dom.canvas.width = Math.round(cssSize * dpr);
dom.canvas.height = Math.round(cssSize * dpr);
ctx.setTransform(dpr, 0, 0, dpr, 0, 0);

view.size = cssSize;
view.cell = cssSize / CONFIG.cols;
	\end{lstlisting}
	\legend{Fonte: \texttt{index.html}, seção~9.}
\end{quadro}

\textbf{Conversão de célula para pixel em um único ponto.} O tamanho da célula é
calculado uma única vez, e todo o desenho deriva dele. Não há coordenada de pixel
literal no código de desenho, exceto as relativas à própria célula (espessura de
borda, deslocamento dos olhos). Essa disciplina é o que permite alterar o número
de colunas da grade sem tocar no código de desenho.

\textbf{Observação de redimensionamento.} O ajuste é reexecutado quando o
contêiner muda de tamanho, por observador de redimensionamento quando disponível
e por evento de janela como alternativa. Isso cobre rotação de tela e
redimensionamento de janela sem recarregar a página.

A Figura~\ref{fig:interface} registra visualmente o resultado da implementação no
estado de fim de jogo: o tema escuro com acentos neon (\texttt{NFR-003}), o
layout centralizado e responsivo (\texttt{NFR-002}), o placar com as duas
grandezas (\texttt{FR-009} e \texttt{FR-010}), a sobreposição de fim de jogo com
o botão de reinício (\texttt{FR-011}) e os quatro controles direcionais
(\texttt{FR-004}).

Cabe uma advertência sobre o estatuto dessa imagem. Ela é \emph{registro visual},
não verificação. Nenhum critério de aceitação é satisfeito ou reprovado com base
em uma captura de tela: os requisitos que ela ilustra são verificados por
observação de comportamento (Tabela~\ref{tab:ca-resultados}) ou, quando não
possuem critério, permanecem sem verificação. Ilustrar um requisito em uma figura
não o promove a verificado.

%% ---------------------------------------------------------------------
%%  FIGURA - Registro visual da interface implementada
%%
%%  A imagem NAO esta em figures/: ela e a captura de tela real do produto,
%%  mantida em images/ na raiz do repositorio (ver \graphicspath em main.tex).
%%  O caminho e resolvido por ../images/screenshot_gaming.png.
%% ---------------------------------------------------------------------
\begin{figure}[htb]
	\caption{Interface implementada, no estado de fim de jogo com placar preenchido}
	\label{fig:interface}
	\centering
	\includegraphics[width=0.74\textwidth]{screenshot_gaming.png}
	\legend{Fonte: arquivo \texttt{images/screenshot\_gaming.png}. Captura de tela do sistema-alvo
	em execução, no estado \texttt{over}. São visíveis: o título da aplicação; o placar com
	pontuação atual e melhor pontuação, ambas em 40; o tabuleiro com o corpo do jogador; a
	sobreposição de fim de jogo com a pontuação final e o botão de reinício; os quatro botões
	direcionais; e a linha de atalhos de teclado. Esta figura é \emph{registro visual} da
	interface, não verificação de critério de aceitação --- a verificação está no
	Capítulo~\ref{cap:resultados}.}
\end{figure}


\section{Acessibilidade e resposta a preferências}
\label{sec:impl-acessibilidade}

Três medidas atendem \texttt{NFR-009} e às diretrizes de acessibilidade para
conteúdo web (\citeonline{w3c2023wcag22}).

\textbf{Preferência de movimento reduzido.} A animação pulsante do alvo é
desativada e as transições de interface são reduzidas quando o sistema operacional
declara essa preferência. A escolha é relevante porque a pulsação é o único
elemento puramente decorativo com movimento contínuo no produto.

\textbf{Foco visível.} Os elementos interativos têm indicação de foco explícita,
o que permite operá-los por teclado. Sem isso, os botões direcionais seriam
inacessíveis a quem não usa apontador.

\textbf{Região viva para leitores de tela.} Mudanças de estado da partida são
anunciadas em região com semântica de estado, de modo que a transição para pausa
ou para fim de jogo não seja silenciosa para tecnologias assistivas.

Registra-se uma limitação: os controles direcionais são botões com rótulo
textual, o que os torna acessíveis, mas a superfície de desenho do jogo não
oferece descrição textual do estado corrente para leitores de tela. Um jogo
acessível a pessoas cegas exigiria uma representação textual do tabuleiro, que
não foi especificada e está fora do escopo deste trabalho.

\section{Verificação de restrições durante a implementação}
\label{sec:impl-restricoes}

A Tabela~\ref{tab:restricoes-verificacao} registra como cada restrição foi
verificada na implementação.

\begin{table}[htb]
	\caption{Verificação das restrições na implementação}
	\label{tab:restricoes-verificacao}
	\centering
	\footnotesize
	\begin{tabularx}{\textwidth}{@{}p{1.7cm} X p{2.6cm}@{}}
		\toprule
		\textbf{ID} & \textbf{Restrição} & \textbf{Verificação} \\
		\midrule
		CON-001 & Artefato único & Inspeção: um único arquivo, um bloco de estilo, um bloco de script \\
		\addlinespace
		CON-002 & Apenas CSS e JavaScript nativos & Inspeção: ausência de referências de rede e de dependências \\
		\addlinespace
		CON-003 & Configuração centralizada & Inspeção: objeto de configuração único, congelado, com sete campos \\
		\addlinespace
		CON-004 & Sem temporizador fixo para simulação & Inspeção: laço por requisição de quadro com acumulador \\
		\bottomrule
	\end{tabularx}
	\legend{Fonte: elaborado pelo autor (2026). A especificação do alvo declara quatro restrições
	(\texttt{CON-001} a \texttt{CON-004}); as restrições de hardware do material de origem do
	framework não têm correspondência no produto (ver P-26).}
\end{table}

\section{Síntese do capítulo}
\label{sec:impl-sintese}

A implementação materializou as sete decisões arquiteturais em 1\,063 linhas
organizadas em treze seções. Os pontos de maior consequência para o
comportamento observável são a ordem entre remoção da cauda e teste de colisão
(ADR-004), a referência usada na validação da fila de entrada (ADR-003) e o
tratamento do estado inicial (FR-012). O Capítulo~\ref{cap:resultados} verifica
se o comportamento resultante corresponde ao especificado.
